
\chapter*{前言}
\addcontentsline{toc}{chapter}{前言}
\markboth{前言}{前言}

本书建立在我们 2011 年出版的《靶向学习：观察性研究与实验研究的因果推断》（Targeted Learning: Causal Inference for Observational and Experimental Studies）之上，是该书的续篇。自这第一部论述因果推断中机器学习方法的著作出版以来，靶向学习取得了多项重要进展。我们认为，有必要出版第二部著作，将这些新成果纳入其中。同时，我们也希望恰当地说明靶向学习方法在蓬勃发展的数据科学领域中所起的作用。

本教材面向统计学、数据科学和公共卫生领域的学者，讨论规模大、结构复杂且动态变化的数据所带来的实际挑战。我们以靶向极大似然估计量（targeted maximum likelihood estimators）的一般构造为模板，提出一条科学路线，把现实世界中的数据科学应用转化为形式化的统计估计问题。这些靶向机器学习算法既能估计所关心的量，又能提供有效的统计推断。在当今时代，数据科学中的靶向学习方法是解决科学问题的重要组成部分。这些技术能够回答复杂问题，例如：如何依据存在时变混杂的纵向数据，确定分配处理的最优规则；又如：如何估计网络等相依数据结构中的其他目标量。《数据科学中的靶向学习》结合软件包和真实数据集给出示范，说明靶向学习对于新一代统计学家和数据科学家的重要意义。

\begin{sourcebox}
\textbf{《数据科学中的靶向学习》的主要特点：}
\begin{enumerate}
\item 提供基于现实数据回答复杂大数据问题所必需的数据分析工具。
\item 介绍既能服务数据科学、又能提供统计推断的机器学习估计量。
\item 入门章节以易于理解的方式，解释复杂纵向数据的因果推断与靶向学习。
\item 包含大量现实应用，展示如何将实际问题转化为统计估计问题，以及如何利用真实数据和靶向统计学习方法，回答所关心的科学问题。
\item 展示在不同实验单位的数据相互依赖时，如何依据实验开展靶向学习，网络描述的依赖结构就是一例。
\item 在保持坚实理论基础的同时，处理规模大、结构复杂且动态变化的数据所带来的实际挑战。
\end{enumerate}
\end{sourcebox}

\section*{全书概要}

与前书相似，《数据科学中的靶向学习》的一个特点，是包含特邀作者撰写的内容，却又不是传统意义上的论文选编。作为本书作者，我们再次投入大量时间修改各章，使符号、风格以及读者熟悉的研究路线保持一致。因此，这第二部靶向学习著作在内容上相互贯通，读来是一部完整统一的教材。

\subsection*{第 I 部分——数据科学中的靶向学习：导论}

第 1 章阐述开展靶向学习的动机，并概述其研究路线：（1）将数据看作随机变量；（2）用统计模型表示对数据实验的真实知识；（3）把科学问题转化为统计目标参数；（4）进行靶向最小损失估计（targeted minimum loss-based estimation，TMLE），并给出相应推断。第 1 章还定义了一个贯穿后文的纵向实例，其灵感来自 HIV 研究中关于“何时开始治疗”的应用问题。

第 2 章回顾纵向数据的因果模型，说明如何用它们定义形式化的因果量或因果查询，从而表达所关心科学问题的答案。我们重点讨论的因果量，是某种干预规则下结局的反事实均值。在序贯随机化假设和正值性假设下，我们给出 g 计算估计目标，将该因果量识别为数据分布的一个函数。至此，统计模型和估计目标均已确定，统计估计问题也随之明确：在因果模型和可识别性假设成立时，这一估计目标等于我们希望得到的因果量；无论这些假设是否成立，它本身都具有有效且值得关注的统计解释。

第 I 部分最后两章专门讨论估计。第 3 章介绍序贯超级学习方法，用于估计反事实结局关于基线协变量的预测函数，这一对象本身就具有独立的研究价值。由于该预测函数的期望等于反事实结局的边际均值，这种序贯学习也为第 4 章定义的 TMLE 提供了初始估计量。第 4 章在贯穿全书的实例中，介绍反事实结局均值的 TMLE。该章通过以下要素，展示计算 TMLE 的一般路线：有效影响曲线；以初始估计量为偏置项、且其损失函数得分的线性张成空间包含有效影响曲线的局部最不利子模型；以及一个迭代算法，每一步都以当前估计量为偏置项，拟合相应最不利子模型的极大似然估计量，并用它更新当前估计量。第 4 章还说明如何利用 TMLE 的影响曲线，对它进行一般分析并开展正式的统计推断。

\subsection*{第 II 部分——核心专题的进一步讨论}

第 II 部分讨论一般 TMLE 的理论和方法进展。在许多估计问题中，TMLE 需要沿着穿过当前更新值的局部最不利子模型，反复最大化对数似然，或者更一般地，最小化损失函数的经验均值，由此迭代更新初始估计量。这种迭代过程可能导致 TMLE 的有限样本表现不稳定，尤其是在数据对目标参数提供的信息十分稀疏时。第 5 章基于所谓通用（典范）最不利子模型，构造一般的一步 TMLE。这类子模型由典范梯度／有效影响曲线唯一刻画。我们还将一步 TMLE 推广到多元目标参数，乃至无限维目标参数。这种 TMLE 的构造思想强烈表明，与迭代 TMLE 相比，一步 TMLE 更稳健、更稳定。本章详细讨论一个生存分析实例，并用模拟说明理论结果。

第 6 章提出一种新的通用估计量，所估计的参数定义为某个损失函数期望的最小化解，例如条件均值或条件密度。我们将其称为\textbf{高度自适应 Lasso（highly adaptive lasso，HAL）估计量}，因为它可以通过如下方式实现：在极高维指示基函数的线性组合中最小化经验风险，同时约束系数绝对值之和不超过某个常数，而该常数本身由交叉验证选取。我们证明，即使在完全非参数模型和高维数据结构下，该估计量也能保证以快于 $n^{-1/4}$ 的速率收敛于真实对象，其中 $n$ 为样本量。

第 7 章证明，只要强正值性假设成立，以 HAL 估计量为初始估计量，或采用候选算法库中包含 HAL 估计量的超级学习器所构造的 TMLE，在很弱的正则条件下就是渐近有效的。TMLE 的构造依赖于有效影响曲线的计算。但在许多估计问题中，有效影响曲线只有隐式表达，依据这一表达直接计算极其困难。第 8 章提出一种基于机器学习的有效影响曲线估计量，从而避免解析计算；使用 HAL 估计量，还可以保证估计精度。尤其是，我们证明，这使得一些以往因有效影响曲线过于复杂而无法构造的 TMLE，变得可以构造。本章通过多种删失数据实例说明这一点，例如区间删失数据和双变量右删失数据。

第 9 章提出一类一般的数据自适应目标参数，使统计学家能够通过数据挖掘来确定感兴趣的目标参数，同时仍获得有效的置信区间。具体而言，我们提出这一类目标参数的交叉验证 TMLE（cross-validated TMLE，CV-TMLE），建立正式的渐近理论，并以连续变量的变量重要性分析展示其应用。随后，第 10 章提出一类一般的协同 TMLE（collaborative TMLE，C-TMLE），在连续统大小的候选估计量集合中，针对目标选择干扰参数估计量。我们证明，与采用专门针对干扰参数本身进行优化的估计量（如超级学习器）来估计干扰参数的 TMLE 相比，该方法在理论上更优。C-TMLE 具有极大的实际重要性，因为它能够保护 TMLE，避免使用会损害其表现的干扰参数拟合结果，例如一个在使用者不知情的情况下纳入了工具变量的倾向评分拟合。

\subsection*{第 III 部分——随机试验}

第 III 部分讨论随机对照试验（randomized controlled trials，RCT）中的 TMLE，包括整群随机对照试验（cluster randomized controlled trials，CRT）。第 11 章构造疫苗位点特异性因果效应的 TMLE：所研究的是，在某个位点与疫苗匹配或不匹配的病毒导致 HIV 感染的时间，如何受到疫苗接种的影响。这种分析能够评价疫苗在不同位点上的有效性，从而为日后的疫苗改进提供方向。TMLE 利用基线协变量和时变协变量来提高效率，并允许存在信息性删失。本章以一项 HIV 疫苗随机对照试验展示该方法。

第 12 章研究样本平均处理效应的 TMLE。样本平均处理效应定义为实际样本中所有个体因果效应的样本均值。本章详细研究稳健统计推断，并以模拟评价所提方法。理论和实践均表明，样本平均处理效应能够比总体平均处理效应估计得更精确。我们还指出，它的另一个优点是，无须将样本视为从某个目标总体中随机抽取的样本。本章以一项比较接受和未接受处理社区的整群随机试验，说明这一点的重要性：这些社区代表东非某个特定地理区域，不能视为随机样本。

第 13 章针对社区数量较少的常见整群随机试验情形，例如只有 30 个社区，提出一种新的数据自适应 TMLE。值得注意的是，这种 TMLE 仍使用超级学习来估计结局回归，其候选库由简单的靶向回归组成，每个回归调整一至两个可能重要的协变量。用于在候选回归估计量中选择最优者的交叉验证准则，以最小化 TMLE 的方差为目标。模拟研究展示了这种数据自适应 TMLE 相对于简单边际估计量在实际表现上的优势。本章的重要意义在于表明：即使样本量很小，TMLE 的核心概念仍然适用。

\subsection*{第 IV 部分——观察性数据}

第 IV 部分讨论利用 TMLE 分析观察性研究。第 14 章针对单时点数据结构 $(W,A,Y)$，构造随机干预因果效应的 TMLE，重点研究依赖于未知处理机制的随机干预。第 15 章和第 16 章展示纵向 TMLE 在复杂观察性纵向研究中的强大应用。第 15 章利用 PROBIT 研究，评价不同母乳喂养方案对儿童发育的因果效应，同时处理整群抽样问题。第 16 章依据 Kaiser Permanente 的大型数据库，评价糖尿病患者采用不同血糖控制强化规则时，对长期事件发生时间结局的影响。第 17 章讨论纵向研究中的因果中介分析。具体而言，该章构造一种新的 TMLE，用于估计单时点处理对生存结局的自然直接或间接因果效应，控制时变中介变量，同时允许存在受时变协变量影响的右删失。

\subsection*{第 V 部分——在线学习}

第 V 部分针对在线时间序列相依数据，构造具有可扩展性的在线超级学习和在线 TMLE。其特例包括人为规定顺序的独立同分布数据；在这种情形下，可以在数据集的不同排序上运行在线估计量。第 18 章提出用于时间序列相依数据的在线超级学习器，并通过建立在线交叉验证选择器的预言不等式，证明其最优性。统计模型假定数据具有 Markov 依赖和平稳性，但不预先指定平稳数据生成机制。第 19 章构造时间序列相依数据的在线 TMLE，估计对序列中某些处理节点进行随机或确定性干预，对未来结局产生的因果影响。这些在线估计量可以用于单条时间序列，此时渐近分析以时间点的数量为基础。在线数据、可扩展性和时间序列相依性都十分重要，因此，本部分提出的模型和方法将推动大量后续研究。特别是在精准医疗时代，能够从单个受试者的实验及数据中学习，具有重要意义。

\subsection*{第 VI 部分——网络}

第 VI 部分讨论由相互连接的单位组成的有限大总体中，随机或确定性干预对平均结局所产生的因果效应及其 TMLE。网络结构随时间被观测，也就是说，对每个单位，我们知道它在各时间点可能依赖的朋友集合。第 20 章介绍纵向网络数据的因果模型和统计模型，并在每个单位都观测到基线协变量 $W$、处理 $A$ 以及不同时间点 $t$ 上的结局 $Y_t$ 时，为所关心的因果效应构造 TMLE。第 21 章关注一个特殊情形：只在一个时间点 $t$ 观测结局，该时间点通常紧随处理 $A$ 之后。在这一特例中，我们证明 TMLE 具有闭式表达，并提供配套 R 软件包。本章详细讨论许多值得关注的问题，并通过模拟评价 TMLE 的实际表现。

第 21 章扩展了关于独立同分布单时点处理数据结构 $(W,A,Y)$ 的大量因果推断文献，允许：（1）一个单位的结局 $Y$ 受到其朋友的处理 $A$ 和基线协变量 $W$ 的影响；（2）一个单位的处理受到其朋友基线协变量的影响。该章还回顾建立 TMLE 渐近正态性的正式渐近理论。需要指出，第 V 部分和第 VI 部分都在样本量为一的问题中构造 TMLE 及其理论：我们只观测到某个复杂实验的一次实现，这个实验可能是单个单位跨越多个时间点的观测，也可能是多个相互依赖的单位在有限时间点上的观测。

\subsection*{第 VII 部分——最优动态规则}

第 VII 部分的三章研究如何估计使平均结局达到最优的动态处理分配规则。各章重点考虑对每个受试者观测基线协变量 $W$、二元处理 $A$ 和结局 $Y$ 的情形，相应推广已在配套论文中给出。第 22 章构造最优规则的超级学习器、最优规则下反事实均值的 TMLE，以及一个数据自适应目标参数的 TMLE；后者定义为估计得到的最优规则下的反事实均值。第 23 章把这项工作扩展到资源约束下的最优规则。

第 24 章提出一种分组序贯自适应随机化设计：根据此前已招募受试者的数据估计最优处理规则，再据此确定下一组受试者的随机化概率。这样，我们既能学习最优规则，又能依据对最优规则的当前最佳估计分配处理。对于这种分组序贯靶向自适应设计，我们提出最优规则下反事实均值的一种新 TMLE。我们证明，只要能够成功学习最优规则，这一 TMLE 就具有渐近一致性和渐近正态性。如果后一个条件不成立，我们仍然能够对数据自适应目标参数进行有效推断，即对估计最优规则下的反事实均值进行推断。也就是说，正如我们总能在随机对照试验中对平均处理效应获得有效推断，对于下面这个复杂得多的因果问题，我们仍保留这一保证：“根据数据学得的这个最优规则估计，其对应的反事实结局均值是多少？”第 24 章解决的问题，是长期以来的多臂老虎机问题的一个低维版本。

\subsection*{第 VIII 部分——专题研究}

第 VIII 部分深入讨论靶向学习领域的一些重要专题。以往靶向学习主要关注路径可微的目标参数。第 25 章研究非路径可微目标参数的估计，例如某一点的密度。我们用由带宽索引的一族路径可微参数逼近目标参数，继而构造 CV-TMLE 及带宽选择器。此外，我们证明，所得 CV-TMLE 以某个未知的自适应速率渐近服从正态分布，该速率取决于底层数据密度的光滑性；并进一步构造渐近有效的置信区间。本章为非路径可微目标参数的靶向学习开辟了一般途径，同时保留正式统计推断的能力。

第 26 章回顾高阶路径可微目标参数的高阶影响函数理论，并说明如何在 TMLE 框架下方便地构造高阶 TMLE。高阶路径可微性的好处在于，可以依据高阶 Taylor 展开构造估计量，这样只须假定更高阶的余项可以忽略，而不必要求二阶余项可以忽略。然而遗憾的是，在符合现实的模型中，多数目标参数只具有一阶路径可微性。如果目标参数仅一阶路径可微，我们证明，基于近似并经过调节的二阶影响函数构造的二阶 TMLE，相对于只针对一阶有效影响函数进行更新的普通 TMLE，可以显著改善有限样本表现。

敏感性分析旨在为估计目标的估计值与所关心因果量之间的差异设定上界。因此，它自然涉及两类偏差：采用有偏估计量所产生的统计偏差，以及违反了从观测数据分布识别因果量所需因果假设而产生的识别偏差。许多敏感性分析因使用有偏估计量，例如参数模型中的回归，而变得难以解释。因此，我们希望采用基于高度自适应超级学习器的 TMLE 等估计量，尽可能保证关于估计目标的统计推断如实可靠。即使实现了这一点，仍然需要针对不可检验假设开展敏感性分析。另一方面，文献提出的许多方法使用参数敏感性模型，导致敏感性参数的解释完全依赖于这些模型的正确性，而这些参数的界限通常被认为可以由外部知识给出。因此，为使敏感性分析透明且有用，必须采用定义明确的敏感性参数。第 27 章提出这样一种非参数敏感性分析方法，使敏感性参数具有非参数意义下的解释，并以一个实际案例，展示该方法的效力和透明性。

当估计量使用机器学习时，非参数自助法通常无法一致地估计其抽样分布；用于现实统计模型的典型 TMLE 就属于这种情形。自助法能够捕捉一阶渐近分析没有涵盖的估计量二阶变异，因此，使自助法也能用于 TMLE 很重要。第 28 章提出靶向自助法，用于估计渐近线性估计量的极限分布。该方法从数据分布 $P_0$ 的 TMLE 估计 $P_n^*$ 中重新抽样，以估计所研究估计量的抽样分布；这里，$P_n^*$ 所针对的目标，是该估计量影响曲线的方差。我们以单时点处理数据结构 $(W,A,Y)$ 下反事实均值的 TMLE，说明这一一般方法。模拟研究评价非参数自助法的失效，以及靶向自助法的性能优势。

第 29 章讨论如何快速计算 TMLE：用基于数据库中规模小得多的受控随机样本所构造的 TMLE，替代基于全部数据的 TMLE，因此得到的估计量不再有效。抽样概率可以是某个对所有个体都可获得、且容易计算的测量量的函数。本章推导使 TMLE 效率最大的最优抽样概率。由此可以考虑分组序贯设计，根据既往数据调整抽样概率，从而使设计最小化 TMLE 的方差。这一一般方法能够把 TMLE 扩展到大型数据集。

最后，第 30 章从历史和哲学角度审视本书主题。在《真理的困境：论统计学、因果性、物理学与科学哲学》一文中，作者讨论数据科学的近期发展对统计学和认识论的一些主要影响。他指出，这些进展带来一种特殊且令人不安的科学观，对统计学和科学哲学而言，这一科学观似乎都不够恰当。随后他说明：鉴于这种反统计、反认识论的立场，改善统计学与科学哲学之间的关系，可以视为一种对自身利益的明智认识，但建立这种联系显然困难重重。如今，对真理的关注似乎至少沿着两条不同的路线展开，认识论与研究方法各司其职，因而形成僵硬而令人遗憾的界线；这种界线也延伸到因果性等其他认识论核心问题。本章详细分析靶向学习在这场争论中的作用和意义，并为数据科学哲学迈出初步步伐。

\subsection*{附录}

若干基础结果的较长证明放在附录中。具体而言，第 A.1 节给出数据自适应目标参数的 CV-TMLE 的一般分析，对应第 9 章。第 A.2 节建立中介分析的三个基础结果，对应第 17 章。第 A.3 节证明时间序列相依数据的在线超级学习器的预言不等式，对应第 18 章。最后，第 A.4 节依据时间序列数据因果目标参数的典范梯度，给出其一阶 Taylor 展开，为第 19 章的 TMLE 分析奠定基础。

\section*{致谢}

本书凝聚了一个庞大团队的共同努力，没有他们，本书便无法完成。我们感谢以下学生和同事对这些方法的贡献与讨论，以及至关重要的协作；其中许多人还为本书撰写了章节：Samrachana Adhikari（Harvard）、Laura B. Balzer（UMass）、Oliver Bembom（Pandora）、David Benkeser（Emory）、Savannah Bergquist（Harvard）、Aurélien Bibaut（Berkeley）、Brian Caffo（Johns Hopkins）、Wilson Cai（Berkeley）、Marco Carone（Washington）、Wilson Cai（Berkeley）、Antoine Chambaz（Université Paris Descartes）、Jeremy Coyle（Berkeley）、Victor De Gruttola（Harvard）、Irina Degtiar（Harvard）、Iván Díaz（Weill Cornell）、Francesa Dominici（Harvard）、Randall Ellis（BU）、Ani Eloyan（Brown）、Bruce Fireman（Kaiser Permanente）、Peter Gilbert（Washington）、Jeff Goldsmith（Columbia）、Susan Gruber（Harvard）、Laura Hatfield（Harvard）、Joe Hogan（Brown）、Alan E. Hubbard（Berkeley）、Nicholas P. Jewell（Berkeley）、Emilien Joly（Université Paris Nanterre）、Cheng Ju（Berkeley）、Chris J. Kennedy（Berkeley）、Tim Layton（Harvard）、Erin LeDell（h2o）、Sam Lendle（Pandora）、Alexander R. Luedtke（Fred Hutch）、Xavier Mary（Université Paris Nanterre）、Alex McDowell（Harvard）、Thomas McGuire（Harvard）、Caleb Miles（Berkeley）、Andrew Mirelman（York）、Ellen Montz（Harvard）、Erica E. M. Moodie（McGill）、Kelly L. Moore（Netflix）、Ian Nason（Harvard）、Romain S. Neugebauer（Kaiser Permanente）、Sharon-Lise Normand（Harvard）、Patrick J. O'Connor（HealthPartners）、Elizabeth L. Ogburn（Johns Hopkins）、Judea Pearl（UCLA）、Maya L. Petersen（Berkeley）、Robert W. Platt（McGill）、Eric C. Polley（Mayo Clinic）、Kristin E. Porter（MDRC）、A. J. Rosellini（BU）、Michael Rosenblum（Johns Hopkins）、Daniel B. Rubin（FDA）、Stephanie Sapp（Google）、Julie A. Schmittdiel（Kaiser Permanente）、Mireille E. Schnitzer（Université de Montréal）、Jasjeet S. Sekhon（Berkeley）、Taki Shinohara（Penn）、Akritee Shrestha（Wayfair）、Michael J. Silverberg（Kaiser Permanente）、Oleg Sofrygin（Berkeley）、Jake Spertus（Harvard）、Richard J. C. M. Starmans（Universiteit Utrecht）、Ori M. Stitelman（Dstillery）、Ira B. Tager（Berkeley）、Boriska Toth（Berkeley）、Catherine Tuglus（Amgen）、Hui Wang（VA）、Yue Wang（Novartis）、C. William Wester（Vanderbilt）、Wenjing Zheng（Netflix）以及 Corwin Zigler（Harvard）。

\begin{flushright}
Mark J. van der Laan\quad 美国加利福尼亚州伯克利\par
Sherri Rose\quad 美国马萨诸塞州波士顿\par
2017 年 3 月
\end{flushright}
